\subsection{Convergent series}

3.1.1. Let \(f = \sum f_{\alpha}\) be a formal series belonging to \(\hat{\mathbf{P}}(\mathbf{E}_1, \ldots, \mathbf{E}_n; \mathbf{F})\) (Appendix, no. A.5). If \(\gamma\) is a continuous semi-norm on \(\mathbf{F}\) and \(\mathbf{R} = (\mathbf{R}_1, \ldots, \mathbf{R}_n)\) a sequence of \(n\) strictly positive real numbers, we set:

\[
\| f \| _ {\gamma , \mathbf {R}} = \sum_ {\alpha} \mathbf {R} ^ {\alpha} \| f _ {\alpha} \| _ {\gamma}.
\]

If \(F\) is a normed space and if \(\gamma\) is the norm of \(F\) , we write \(\| f \|_{\mathbb{R}}\) instead of \(\| f \|_{\gamma, \mathbb{R}}\) .

The set \(\mathcal{H}_{\mathbb{R}}(\mathbf{E}_{1},\ldots,\mathbf{E}_{n};\mathbf{F})\) of \(f\in\hat{\mathbf{P}}(\mathbf{E}_{1},\ldots,\mathbf{E}_{n};\mathbf{F})\) such that \(\|f\|_{\gamma,\mathbb{R}}\) is finite for every continuous seminorm \(\gamma\) on \(\mathbf{F}\) is a vector subspace of \(\hat{\mathbf{P}}(\mathbf{E}_{1},\ldots,\mathbf{E}_{n};\mathbf{F})\) . One has

\[
\mathcal {H} _ {\mathbb {R}} (\mathrm{E} _ {1}, \dots , \mathrm{E} _ {n}; \mathrm{F}) \subset \mathcal {H} _ {\mathbb {R} ^ {\prime}} (\mathrm{E} _ {1}, \dots , \mathrm{E} _ {n}; \mathrm{F})
\]

whenever \(R_{i} \geqslant R_{i}'\) for \(1 \leqslant i \leqslant n\) . The union of the

\[
\mathcal {H} _ {\mathbf {R}} (\mathrm{E} _ {1}, \dots , \mathrm{E} _ {n}; \mathrm{F})
\]

is a vector subspace, denoted \(\mathcal{H}(E_{1},\ldots,E_{n};F)\) , of \(\hat{\mathrm{P}}(E_{1},\ldots,E_{n};F)\) , whose elements are called convergent series on the product of the \(E_{i}\) , with values in \(F\) . It depends only on the topology of the \(E_{i}\) and not on their norms. One can therefore speak of the space \(\mathcal{H}(E_{1},\ldots,E_{n};F)\) when the \(E_{i}\) are normable spaces, without having to choose a norm on each \(E_{i}\) .

3.1.2. The map \(f\mapsto\|f\|_{\gamma,\mathbb{R}}\) is a seminorm on \(\mathcal{H}_{\mathbb{R}}(\mathrm{E}_{1},\ldots,\mathrm{E}_{n};\mathrm{F})\) . The topology defined by these seminorms as \(\gamma\) runs over the set of continuous seminorms on F (or simply a set of seminorms defining the topology of F) is Hausdorff. If F is normable (resp. complete), so is \(\mathcal{H}_{\mathbb{R}}(\mathrm{E}_{1},\ldots,\mathrm{E}_{n};\mathrm{F})\) . The seminorms on \(\mathcal{H}(\mathrm{E}_{1},\ldots,\mathrm{E}_{n};\mathrm{F})\) whose restriction to each \(\mathcal{H}_{\mathbb{R}}\) is continuous define a Hausdorff topology on \(\mathcal{H}(\mathrm{E}_{1},\ldots,\mathrm{E}_{n};\mathrm{F})\) , which depends only on the topologies of the \(\mathrm{E}_{i}\) . The injection of \(\mathcal{H}(\mathrm{E}_{1},\ldots,\mathrm{E}_{n};\mathrm{F})\) into \(\hat{\mathrm{P}}(\mathrm{E}_{1},\ldots,\mathrm{E}_{n};\mathrm{F})\) is continuous.

3.1.3. The canonical isomorphism of \(\hat{\mathbf{P}}(\mathbf{E};\mathbf{F})\) onto \(\hat{\mathbf{P}}(\mathbf{E}_{1},\ldots,\mathbf{E}_{n};\mathbf{F})\) gives by restriction an isomorphism of topological vector spaces from \(\mathcal{H}(\mathbf{E};\mathbf{F})\) onto \(\mathcal{H}(\mathbf{E}_{1},\ldots,\mathbf{E}_{n};\mathbf{F})\) .

3.1.4. Let \(f \in \mathcal{H}(\mathrm{E}_1, \ldots, \mathrm{E}_n; \mathrm{F})\) and let \(J(f)\) be the set of \(R \in (\mathbf{R}_+^*)^n\) such that \(f \in \mathcal{H}_{\mathbb{R}}(\mathrm{E}_1, \ldots, \mathrm{E}_n; \mathrm{F})\) . The interior \(I(f)\) of \(J(f)\) is called the strict convergence indicator of \(f\) . It consists of the set of \(R\) for which there exists an \(R' \in J(f)\) with \(0 < R_i < R'_i\) for \(1 \leqslant i \leqslant n\) . We denote by \(\Omega(f)\) the set of points \((\log R_1, \ldots, \log R_n)\) of \(R^n\) for \(R \in I(f)\) : it is a convex subset of \(R^n\) .

When \(n=1\) , the set \(\mathbf{I}(f)\) is an interval \([0,\rho(f)]\) of \(\mathbf{R}\) and \(\rho(f)\) is called the strict radius of convergence of \(f\) . It is also the supremum (finite or \(+\infty\) ) of the set of real numbers \(R>0\) such that for every continuous seminorm \(\gamma\) on \(F\) , there exists a constant \(M\) such that \(\|f_{m}\|_{\gamma}\leqslant MR^{-m}\) (with \(f=\sum f_{m},f_{m}\in P_{m}(E;F)\) ) for every integer \(m\geqslant0\) .

The set of points \(x = (x_i) \in \mathbf{E}_1 \times \cdots \times \mathbf{E}_n\) such that there exists \(\mathbf{R} \in \mathbf{I}(f)\) with \(\|x_i\| \leqslant R_i\) for all i, is called the strict domain of convergence of f and is denoted \(\mathbf{C}(f)\) . It is also the interior of the set of points x for which there exists \(\mathbf{R} \in \mathbf{J}(f)\) with \(\|x_i\| < R_i\) for all i.

3.1.5. For \(f_{\alpha} \in \mathbf{P}_{\alpha}(\mathbf{E}_1, \ldots, \mathbf{E}_n; \mathbf{F})\) and for every continuous seminorm \(\gamma\) on \(\mathbf{F}\) , set:

\[
\| f _ {\alpha} \| _ {\gamma} = \sup _ {\| x _ {i} \| \leqslant 1} \| f _ {\alpha} (x _ {1}, \dots , x _ {n}) \| _ {\gamma}.
\]

We then have the inequalities:

\[
\left\| \left| f _ {\alpha} \right| \right\| _ {\gamma} \leqslant \left\| f _ {\alpha} \right\| _ {\gamma} \leqslant \frac {\alpha^ {\alpha}}{\alpha !} \left\| \left| f _ {\alpha} \right| \right\| _ {\gamma}.
\]

For \(f = \sum f_{\alpha} \in \hat{\mathbf{P}}(\mathrm{E}_1, \ldots, \mathrm{E}_n; \mathbf{F})\) and \(\mathbf{R} \in (\mathbf{R}_{+}^{*})^n\) , set:

\[
\| | f \| | _ {\gamma , \mathbf {R}} = \sum_ {\alpha} \mathbf {R} ^ {\alpha} \| | f _ {\alpha} \| | _ {\gamma}
\]

and let us denote by \(\tilde{\mathcal{H}}_{\mathbb{R}}(\mathrm{E}_{1},\ldots ,\mathrm{E}_{n};\mathrm{F})\) the vector subspace of

\[
\hat {\mathbf {P}} (\mathbf {E} _ {1}, \dots , \mathbf {E} _ {n}; \mathbf {F})
\]

formed by the \(f\) such that \(\| | f \| | _ {\gamma,\mathbb{R}}\) is finite for every \(\gamma\) , equipped with the topology defined by the seminorms \(f\mapsto\| | f \| | _ {\gamma,\mathbb{R}}\) . One has \(\mathcal{H}_{\mathbb{R}}\subset\tilde{\mathcal{H}}_{\mathbb{R}}\) and the injection of \(\mathcal{H}_{\mathbb{R}}\) into \(\tilde{\mathcal{H}}_{\mathbb{R}}\) is continuous. The space \(\mathcal{H}(\mathrm{E}_{1},\ldots,\mathrm{E}_{n};\mathrm{F})\) is the union of the \(\tilde{\mathcal{H}}_{\mathbb{R}}(\mathrm{E}_{1},\ldots,\mathrm{E}_{n};\mathrm{F})\), and its topology is the finest locally convex topology making the injections of \(\mathcal{H}_{\mathbb{R}}\) into \(\mathcal{H}\) continuous.

If \(f \in \mathcal{H}(E_1, \ldots, E_n; F)\), one calls the convergence indicator the interior \(\tilde{I}(f)\) of the set \(\tilde{J}(f)\) of \(R \in (\mathbf{R}_+^*)\) such that \(f \in \mathcal{H}_{\mathbb{R}}\). We have:

\[
e ^ {- 1} \mathbf {I} (f) \subset \tilde{\mathbf {I}} (f) \subset \mathbf {I} (f)
\]

(where e is the base of natural logarithms). Starting from \(\tilde{I}(f)\) , one defines as in 3.1.4. the domain of convergence \(\tilde{C}(f)\) and, when \(n = 1\) , the radius of convergence \(\tilde{\rho}(f)\) . In particular, we have:

\[
e ^ {- 1} \tilde {\rho} (f) \leqslant \rho (f) \leqslant \tilde {\rho} (f).
\]

3.1.6. For \(\mathbf{R} \in (\mathbf{R}_{+}^{*})^{n}\), we call the (closed) polyball centered at 0 and with radius \(\mathbf{R}\) in E the set \(\mathbf{B}(\mathbf{R})\) of \(x \in \mathbf{E}\) such that \(\|x_{i}\| \leqslant \mathbf{R}_{i}\) for all \(i\). If \(\dim \mathbf{E}_{i} = 1\), we also say polydisk. If \(f \in \mathcal{H}(\mathbf{E}_{1}, \ldots, \mathbf{E}_{n}; \mathbf{F})\), the domain of convergence (resp. strict) of \(f\) is the union of the polyballs \(\mathbf{B}(\mathbf{R})\) for \(\mathbf{R} \in \tilde{\mathbf{I}}(f)\) (resp. \(\mathbf{R} \in \mathbf{I}(f)\)).

3.1.7. Let \(f \in \mathcal{H}(E_1, \ldots, E_n; F)\) . Suppose F is quasi-complete. For every \(x \in \tilde{\mathbf{C}}(f)\) , the family of \(f_{\alpha}(x)\) is summable in F. Its sum, denoted \(\hat{f}(x)\) or simply \(f(x)\) , is a continuous function on \(\tilde{\mathbf{C}}(f)\) . More precisely, for every R such that \(f \in \tilde{\mathcal{H}}_{\mathbb{R}}\) , the family of \(f_{\alpha}(x)\) is uniformly summable for \(x \in B(R)\) . The map \(f \mapsto \hat{f}\) is a continuous linear injective map from \(\tilde{\mathcal{H}}_{\mathbb{R}}\) into the space of bounded continuous functions on \(B(R)\) , equipped with the topology of uniform convergence.

3.1.8. Let \(F_1, \ldots, F_m\) be separated polynormed spaces and let \(u\) be a continuous \(m\) -linear map from \(F_1 \times \cdots \times F_m\) to \(F\) . Let \(f_i \in \mathcal{H}(E_1, \ldots, E_n; F_i)\) for \(1 \leqslant i \leqslant m\) . The formal series \(u(f_1, \ldots, f_m)\) belongs to \(\mathcal{H}(E_1, \ldots, E_n; F)\) . We have:

\[
\mathrm{C} (u (f _ {1}, \dots , f _ {m})) \supset \bigcap_ {i} \mathrm{C} (f _ {i})
\]

\[
\tilde {C} (u (f _ {1}, \dots , f _ {m})) \supset \bigcap_ {i} \tilde {C} (f _ {i})
\]

and if \(x \in \bigcap_{i} \tilde{C}(f_i)\) , F and the \(F_i\) being quasi-complete, we have:

\[
u (f _ {1}, \dots , f _ {m}) (x) = u (f _ {1} (x), \dots , f _ {m} (x)).
\]

3.1.9. Let \(\mathbf{F}_1, \ldots, \mathbf{F}_m\) be complete normed spaces and suppose \(\mathbf{F}\) quasi-complete. Let \(\mathbf{f} = (f_i)_{1 \leqslant i \leqslant m}\) with \(f_i \in \mathcal{H}(\mathbf{E}_1, \ldots, \mathbf{E}_n; \mathbf{F}_i)\) and \(g \in \mathcal{H}(\mathbf{F}_1, \ldots, \mathbf{F}_m; \mathbf{F})\) , such that \((f_i(0))_{1 \leqslant i \leqslant m}\) belongs to the strict domain of convergence of \(g\) . Then, for every \(\alpha \in \mathbb{N}^m\) , the formal series \(g_\alpha \circ \mathbf{f}\) belongs to \(\mathcal{H}(\mathbf{E}_1, \ldots, \mathbf{E}_n; \mathbf{F})\) and the family of \(g_\alpha \circ \mathbf{f}\) is summable in \(\mathcal{H}(\mathbf{E}_1, \ldots, \mathbf{E}_n; \mathbf{F})\) and a fortiori in \(\hat{\mathbf{P}}(\mathbf{E}_1, \ldots, \mathbf{E}_n; \mathbf{F})\) . Its sum will be denoted \(g \circ \mathbf{f}\) .

More precisely, there exists \(\mathbb{R}\in\bigcap_{i}\mathrm{I}(f_{i})\) and \(\mathbb{R}^{\prime}\in\mathrm{I}(g)\) such that \(\|f_{i}\|_{\mathbb{R}}<R_{i}^{\prime}\) for \(1\leqslant i\leqslant m\) . Under these conditions, the formal series \(g_{\alpha}\circ f\) belongs to \(\mathcal{H}_{\mathbb{R}}(\mathrm{E}_{1},\ldots,\mathrm{E}_{n};\mathrm{F})\) and the family of \(g_{\alpha}\circ f\) is summable in \(\mathcal{H}_{\mathbb{R}}(\mathrm{E}_{1},\ldots,\mathrm{E}_{n};\mathrm{F})\) . Finally, if \(x\in\mathrm{B}(\mathbb{R})\) , then \(\mathbf{f}(x)=(f_{i}(x))\) belongs to \(\mathrm{B}(\mathrm{R}^{\prime})\subset\mathrm{F}_{1}\times\cdots\times\mathrm{F}_{m}\) and we have:

\[
g (\mathbf {f} (x)) = (g \circ \mathbf {f}) (x).
\]

3.1.10. Suppose that \(E_{i} = K\) for \(1 \leqslant i \leqslant n\) . The space \(\hat{\mathbf{P}}(\mathbf{K}^{n}; \mathbf{F})\) is then identified with the space of formal power series in n indeterminates \(X_{1}, \ldots, X_{n}\) with coefficients in F and an element of \(\hat{\mathbf{P}}(\mathbf{K}^{n};\mathbf{F})\) is written:

\[
f = \sum_ {\alpha} X ^ {\alpha} c _ {\alpha} \quad \text {   with   } c _ {\alpha} \in F.
\]

If \(\mathbf{R} \in (\mathbf{R}_{+}^{*})^{n}\) and if \(\gamma\) is a continuous semi-norm on \(\mathbf{F}\) , we have:

\[
\| | f \| | _ {\gamma , \mathbf {R}} = \| f \| _ {\gamma , \mathbf {R}} = \sum_ {\alpha} \mathbf {R} ^ {\alpha} \| c _ {\alpha} \| _ {\gamma}
\]

One has \(\mathbf{I}(f) = \tilde{\mathbf{I}}(f)\) , \(\mathbf{C}(f) = \tilde{\mathbf{C}}(f)\) and, when \(n = 1\) , \(\rho(f) = \tilde{\rho}(f)\) .

The space \(\mathcal{H}(\mathbf{K}^{n};\mathbf{K})\) of convergent series with coefficients in \(\mathbf{K}\) is also denoted \(\mathbf{K}\{\{\mathbf{X}_{1},\ldots ,\mathbf{X}_{n}\}\}\) ; it is a subalgebra of \(\mathbf{K}[[\mathbf{X}_1,\dots ,\mathbf{X}_n]]\) . The space \(\mathcal{H}(\mathbf{K}^n;\mathbf{F})\) is a module over \(\mathbf{K}\{\{\mathbf{X}_1,\ldots ,\mathbf{X}_n\}\}\) and if \(\mathbf{F}\) is finite-dimensional, this module is identified with \(\mathbf{K}\{\{\mathbf{X}_1,\ldots ,\mathbf{X}_n\}\} \otimes_{\mathbf{K}}\mathbf{F}\) .

3.1.11. Let \(\mathbf{f} \in \mathcal{H}(\mathbf{K}^n; \mathbf{K}^m)\), represented by a system of \(m\) convergent series \(f_j(X_1, \ldots, X_n)\), with coefficients in K. Let likewise \(\mathbf{g} \in \mathcal{H}(\mathbf{K}^m; \mathbf{K}^p)\), represented by a system of \(p\) convergent series \(g_k(Y_1, \ldots, Y_m) = \sum g_{k,\beta} Y^\beta\). The element \(\mathbf{h} = \mathbf{g} \circ \mathbf{f}\) of \(\mathcal{H}(\mathbf{K}^n; \mathbf{K}^p)\) (cf. 3.1.9) is represented by \(p\) formal series \(h_k(X_1, \ldots, X_n)\) determined as follows: for \(\alpha \in \mathbf{N}^n\) and \(\beta \in \mathbf{N}^m\), let \(c_{\alpha,\beta}\) be the coefficient of \(X^\alpha\) in the formal series \(\mathbf{f}^\beta = \prod f_j^{\beta_j}\); then the family \((g_{k,\beta} c_{\alpha,\beta})_{\beta \in \mathbb{N}^m}\) is summable in K and has as its sum the coefficient of \(X^\alpha\) in \(h_k\).

3.1.12. Suppose F is quasi-complete and let \(\hat{\mathbf{E}}_{i}\) be the completion of \(\mathbf{E}_{i}\). A continuous polynomial on \(\mathbf{E}_{1} \times \cdots \times \mathbf{E}_{n}\) taking values in F extends by continuity to a polynomial-continuous map on \(\hat{\mathbf{E}}_{1} \times \cdots \times \hat{\mathbf{E}}_{n}\) with values in F. From this we deduce a bijection \(j\) of \(\hat{\mathbf{P}}(\mathbf{E}_{1}, \ldots, \mathbf{E}_{n}; \mathbf{F})\) onto \(\hat{\mathbf{P}}(\hat{\mathbf{E}}_{1}, \ldots, \hat{\mathbf{E}}_{n}; \mathbf{F})\) . If \(f \in \mathcal{H}(\mathbf{E}_{1}, \ldots, \mathbf{E}_{n}; \mathbf{F})\) , then \(j(f) \in \mathcal{H}(\hat{\mathbf{E}}_{1}, \ldots, \hat{\mathbf{E}}_{n}; \mathbf{F})\) and conversely. The strict convergence indicators of \(f\) and of \(j(f)\) are the same.

3.1.13. Suppose that K = R, but that F is equipped with a complex vector space structure compatible with its real vector space structure. Let \(E_{i}^{C} = E_{i} \otimes_{R} C\) . If \(y \in E_{i}^{C}\) , set:

\[
\| y \| = \inf \sum_ {k} | a _ {k} |. \| x _ {k} \|,
\]

the infimum being taken over all finite families of pairs \((x_{k}, a_{k}) \in \mathrm{E}_{i} \times \mathbf{C}\) such that \(y = \sum_{k} x_{k} \otimes a_{k}\) . One thus obtains a norm on the complex vector space \(E_{i}^{c}\) ; this norm extends the given norm on \(E_{i}\) if one agrees to identify \(x \in E_{i}\) with \(x \otimes 1\) . Let \(h\) be an R-polynomial-continuous map on \(E_{1} \times \cdots \times E_{n}\), taking values in F, and homogeneous of multidegree \(\alpha\); there then exists a C-polynomial-continuous \(\tilde{h}\) on \(E_{1}^{c} \times \cdots \times E_{n}^{c}\), and a unique one, with values in F, extending h, and homogeneous of multidegree \(\alpha\). One has

\[
\| \tilde {h} \| _ {\gamma} = \| h \| _ {\gamma}
\]

for every continuous seminorm γ on the complex vector space F.

\[
f = \sum_ {\alpha} f _ {\alpha} \in \mathcal {H} (\mathrm{E} _ {1}, \dots , \mathrm{E} _ {n}; \mathrm{F}), \text {   then   } \tilde {f} = \sum_ {\alpha} \tilde {f} _ {\alpha} \in \mathcal {H} (\mathrm{E} _ {1} ^ {\mathrm{c}}, \dots , \mathrm{E} _ {n} ^ {\mathrm{c}}; \mathrm{F}).
\]

The series \(f\) and \(\tilde{f}\) have the same strict convergence indicator (and the same strict radius of convergence when \(n = 1\) ).

Conversely, suppose K = C. Let \(E_{i}^{0}\) and \(F^{0}\) the spaces over R obtained by restriction of scalars. If \(f_{\alpha} \in \mathrm{P}_{\alpha}(\mathrm{E}_{1}, \ldots, \mathrm{E}_{n}; \mathrm{F})\) , then \(f_{\alpha} \in \mathrm{P}_{\alpha}(\mathrm{E}_{1}^{0}, \ldots, \mathrm{E}_{n}^{0}; \mathrm{F}^{0})\) . If \(f = \sum_{\alpha} f_{\alpha} \in \mathcal{H}(\mathrm{E}_{1}, \ldots, \mathrm{E}_{n}; \mathrm{F})\) then the formal series \(f^{0} = \sum_{\alpha} f_{\alpha} \in \hat{\mathrm{P}}(\mathrm{E}_{1}^{0}, \ldots, \mathrm{E}_{n}^{0}; \mathrm{F}^{0})\) is a convergent series. The convergence indicators (resp. strict convergence indicators) of f and \(f^{0}\) are identical, and we have \(f(x) = f^{0}(x)\) for all \(x \in \tilde{\mathbf{C}}(f) = \tilde{\mathbf{C}}(f^{0})\) .

